% Chapter 07: Publishing
% Packaging, versioning, and distributing Domain Kits

\chapter{Publishing}
\label{ch:publishing}

% ============================================================
\section{Local-First Model}
\label{sec:local-first}
% ============================================================

Gert Domain Kits follow a \textbf{local-first} distribution model. There is no
central Kit registry, no cloud service, and no telemetry. A Kit is a directory
containing a manifest and binaries. It can be:

\begin{itemize}
  \item Checked into a Git repository alongside runbooks.
  \item Distributed as a tarball (\texttt{.tar.gz}).
  \item Installed from a package manager (\texttt{apt}, \texttt{brew},
    \texttt{npm}, etc., if packaged).
  \item Symlinked from a shared filesystem.
  \item Vendored into a project directory.
\end{itemize}

This aligns with gert's local-first philosophy: runbooks and their dependencies
should be portable, versionable, and executable without network access.

% ============================================================
\section{Referencing a Kit from a Runbook}
\label{sec:kit-reference}
% ============================================================

Runbooks reference a Kit in two places:

\paragraph{1. In the \texttt{apiVersion} field:}

\begin{minted}{yaml}
apiVersion: runbook/v1+compliance/v1
\end{minted}

The suffix (\texttt{compliance/v1}) tells gert which Kit schema to load. This
is declared in the Kit manifest under \texttt{schema.apiVersion-suffix}.

\paragraph{2. In the \texttt{meta.kit} field (optional but recommended):}

\begin{minted}{yaml}
meta:
  title: "Quarterly access review"
  kit:
    name: gert.compliance       # Kit identifier
    version: "^1.2.0"            # Semver range
\end{minted}

The \texttt{meta.kit} field allows explicit version pinning. If the installed
Kit version does not satisfy the constraint, \texttt{gert validate} rejects the
runbook.

% ============================================================
\section{Kit Discovery}
\label{sec:kit-discovery}
% ============================================================

When gert encounters a runbook with \texttt{apiVersion: runbook/v1+<suffix>},
it searches for the corresponding Kit in the following order:

\begin{enumerate}
  \item \textbf{Project-local Kits:} \texttt{./kits/<kit-name>/}
  \item \textbf{User Kits:} \texttt{\textasciitilde/.gert/kits/<kit-name>/}
  \item \textbf{\texttt{GERT\_KIT\_PATH} directories:} Each directory in the
    colon-separated \texttt{GERT\_KIT\_PATH} environment variable (similar to
    \texttt{PATH}).
  \item \textbf{System Kits:} \texttt{/usr/local/share/gert/kits/<kit-name>/}
    (platform-specific).
\end{enumerate}

\paragraph{Resolution algorithm:}

\begin{enumerate}
  \item Extract the Kit identifier from \texttt{apiVersion-suffix}. Example:
    \texttt{compliance/v1} → look for a Kit whose
    \texttt{schema.apiVersion-suffix} is \texttt{compliance/v1}.
  \item Search each directory for a \texttt{gert-kit.yaml} manifest.
  \item Load the manifest and check:
    \begin{itemize}
      \item Does \texttt{schema.apiVersion-suffix} match?
      \item If \texttt{meta.kit.version} is specified in the runbook, does the
        Kit version satisfy the constraint?
      \item Is the Kit compatible with this gert core version
        (\texttt{compatibility.gert-core})?
    \end{itemize}
  \item If a match is found, load the Kit. Otherwise, emit an error: ``Kit not
    found: compliance/v1''.
\end{enumerate}

% ============================================================
\section{Packaging}
\label{sec:packaging}
% ============================================================

A Kit is a directory. The simplest packaging is a tarball:

\begin{verbatim}
tar -czf gert.compliance-1.2.0.tar.gz gert.compliance/
\end{verbatim}

Users extract it:

\begin{verbatim}
tar -xzf gert.compliance-1.2.0.tar.gz -C ~/.gert/kits/
\end{verbatim}

Gert will discover the Kit at \texttt{\textasciitilde/.gert/kits/gert.compliance/}.

\paragraph{Platform-specific binaries:}

If your Kit compiler is written in Go, compile binaries for multiple platforms:

\begin{verbatim}
GOOS=linux GOARCH=amd64 go build -o bin/gert-kit-compliance-linux-amd64 ./cmd/compiler
GOOS=darwin GOARCH=arm64 go build -o bin/gert-kit-compliance-darwin-arm64 ./cmd/compiler
GOOS=windows GOARCH=amd64 go build -o bin/gert-kit-compliance-windows-amd64.exe ./cmd/compiler
\end{verbatim}

Ship all binaries in the Kit directory. Gert selects the appropriate binary at
runtime based on \texttt{GOOS}/\texttt{GOARCH}.

\paragraph{Alternative: Universal binaries with \texttt{gert kit install}:}

If gert provides a \texttt{gert kit install} command in the future, it could:

\begin{verbatim}
gert kit install gert.compliance@1.2.0 --from https://github.com/acme/gert-kit-compliance/releases/download/v1.2.0/gert.compliance-1.2.0.tar.gz
\end{verbatim}

This is future work. For now, users install Kits manually.

% ============================================================
\section{Versioning and Pinning}
\label{sec:versioning-pinning}
% ============================================================

\paragraph{Semantic versioning:}

Kit versions MUST follow semver:

\begin{description}
  \item[\texttt{1.0.0}] --- Initial stable release.
  \item[\texttt{1.1.0}] --- Added new step type (backward-compatible).
  \item[\texttt{1.2.0}] --- Added optional field to existing step type
    (backward-compatible).
  \item[\texttt{2.0.0}] --- Renamed required field (breaking change).
\end{description}

\paragraph{Pinning in runbooks:}

Exact pin:

\begin{minted}{yaml}
meta:
  kit:
    name: gert.compliance
    version: "1.2.0"         # Must be exactly 1.2.0
\end{minted}

Range pin (recommended):

\begin{minted}{yaml}
meta:
  kit:
    name: gert.compliance
    version: "^1.2.0"        # Any 1.x >= 1.2.0
\end{minted}

The caret (\texttt{\^}) allows minor and patch updates but blocks major updates.

\paragraph{Validation:}

\texttt{gert validate} checks that the installed Kit version satisfies the pin:

\begin{verbatim}
$ gert validate runbook.yaml
Error: Kit version mismatch
  Required: gert.compliance ^1.2.0
  Installed: gert.compliance 2.0.0
  Location: ~/.gert/kits/gert.compliance/
\end{verbatim}

% ============================================================
\section{Distribution Options}
\label{sec:distribution-options}
% ============================================================

\paragraph{Option 1: Git submodule}

Vendor the Kit into your runbook repository:

\begin{verbatim}
git submodule add https://github.com/acme/gert-kit-compliance kits/gert.compliance
git submodule update --init
\end{verbatim}

Gert will discover the Kit at \texttt{./kits/gert.compliance/}.

\paragraph{Option 2: Local path}

Copy the Kit directory into your project:

\begin{verbatim}
cp -r ~/Downloads/gert.compliance-1.2.0 ./kits/gert.compliance
\end{verbatim}

\paragraph{Option 3: Archived bundle}

Distribute a tarball via GitHub Releases, a web server, or internal artifact
repository:

\begin{verbatim}
curl -L https://github.com/acme/gert-kit-compliance/releases/download/v1.2.0/gert.compliance-1.2.0.tar.gz | \
  tar -xz -C ~/.gert/kits/
\end{verbatim}

\paragraph{Option 4: Package manager}

For broader distribution, package the Kit for \texttt{apt}, \texttt{brew},
\texttt{npm}, etc.

Example \texttt{brew} formula:

\begin{verbatim}
class GertKitCompliance < Formula
  desc "Compliance and audit-focused Domain Kit for gert"
  homepage "https://github.com/acme/gert-kit-compliance"
  url "https://github.com/acme/gert-kit-compliance/archive/v1.2.0.tar.gz"
  sha256 "..."

  def install
    bin.install "bin/gert-kit-compliance"
    prefix.install "gert-kit.yaml", "schemas"
  end
end
\end{verbatim}

Users install with:

\begin{verbatim}
brew install gert-kit-compliance
\end{verbatim}

The formula installs binaries to \texttt{/usr/local/bin/} and Kit artifacts to
\texttt{/usr/local/share/gert/kits/gert.compliance/}.

% ============================================================
\section{Updating a Kit}
\label{sec:updating-kit}
% ============================================================

\paragraph{For runbook authors:}

\begin{enumerate}
  \item Check the Kit changelog to see what changed.
  \item Update the \texttt{meta.kit.version} pin in your runbook.
  \item Run \texttt{gert validate} to ensure the runbook is still valid.
  \item If the Kit has breaking changes (major version bump), update your
    runbook syntax to match the new schema.
\end{enumerate}

\paragraph{For Kit developers:}

\begin{enumerate}
  \item Bump the version in \texttt{gert-kit.yaml}.
  \item Update the changelog.
  \item Tag the release in Git: \texttt{git tag v1.3.0 \&\& git push --tags}.
  \item Build and package binaries for all platforms.
  \item Publish the tarball to GitHub Releases or your distribution mechanism.
\end{enumerate}

% ============================================================
\section{Example: Publishing Checklist}
\label{sec:publishing-checklist}
% ============================================================

\begin{enumerate}
  \item \textbf{Test:} Run \texttt{make test} to ensure all fixtures pass.
  \item \textbf{Version bump:} Update \texttt{meta.version} in \texttt{gert-kit.yaml}.
  \item \textbf{Changelog:} Document changes in \texttt{CHANGELOG.md}.
  \item \textbf{Build binaries:} Compile for Linux, macOS, Windows.
  \item \textbf{Tag release:} \texttt{git tag v1.2.0 \&\& git push --tags}.
  \item \textbf{Package:} \texttt{tar -czf gert.compliance-1.2.0.tar.gz gert.compliance/}
  \item \textbf{Publish:} Upload tarball to GitHub Releases.
  \item \textbf{Announce:} Notify users via documentation, blog post, or release notes.
\end{enumerate}
